\clearpage
\fancyhead[L]{\small BESSEL COMET VON 1807.}
\hypertarget{native-535lower}{}\label{native-535lower}
\begin{center}\small Göttingische gelehrte Anzeigen. Stück 26. Seite 250..256. 1811 Februar 16.\\[3pt]\rule{.72\linewidth}{.3pt}\end{center}

\emph{Untersuchungen über die scheinbare und wahre Bahn des im Jahre \textup{1807} erschienenen grossen Cometen, von} F. W. \textsc{Bessel}, Professor der Astronomie in Königsberg. Bei \textsc{Friedr. Nicolovius}. Königsberg. 82 Seiten in Quart.

Der Comet von 1807, glänzender, als seit vielen Jahren einer sichtbar gewesen war, galt bekanntlich für eines der merkwürdigsten himmlischen Phänomene; noch interessanter aber war er den Astronomen wegen seiner langen Sichtbarkeit, welche die Beobachter gleichsam wetteifernd zu einer recht vollständigen und scharfen Bestimmung seiner scheinbaren Bewegung benutzten. Auch die Berechnung seiner parabolischen Bahn beschäftigte mehrere Astronomen: Niemand hat indessen auf die Untersuchung der Theorie dieses Cometen so viele Sorgfalt verwandt, als der Verfasser der vorliegenden Schrift, welcher bekanntlich sich auch schon um verschiedene ältere Cometen, besonders den von 1769, sehr verdient gemacht hatte. Ein grosser Theil der Untersuchungen des Verfassers über den Cometen von 1807 war schon durch die Monatliche Crrespondenz und das Astronomische Jahrbuch bekannt gemacht; gegenwärtige Schrift hat zum Zwecke, alles dieses im Zusammenhange, und zugleich die spätern Arbeiten vorzutragen, wodurch er die letzte Hand an die Bestimmung der Theorie dieses Cometen legte. Diese Schrift gibt einen neuen Beweis sowohl von dem jetzigen Zustande der Beobachtungskunst, als von der Kraft des Calculs, wenn eine so geschickte Hand ihn leitet.

Der \emph{erste Abschnitt} beschäftigt sich mit der scheinbaren Bahn des Cometen am Himmel. Es war nicht Hrn. \textsc{Bessels} Absicht, alle Beobachtungen des Cometen zu sammeln, sondern nur so viele mit Sorgfalt angestellte zu vereinigen, als zu einer hinreichenden Darstellung der scheinbaren Bewegung des Cometen während seiner ganzen Sichtbarkeit vom 22. September 1807 bis 27. März 1808 erforderlich waren. Er theilt daher zuvörderst seine eigenen, in Lilienthal mit dem Kreismicrometer gemachten, Beobachtungen mit, die bei der gewissenhaften Sorgfalt und der hierin ausgezeichneten Geschicklichkeit des Verfassers ein besonderes Vertrauen verdienen. Verschiedene bei dieser Gelegenheit mitgetheilte Vorsichtsregeln und Bemerkungen über diese Gattung von Beobachtungen werden vielen Lesern sehr willkommen sein. Den Durchmesser des Gesichtsfeldes bestimmte \textsc{Bessel} durch solche Sternpaare, deren Declinationen genau bekannt und beinahe um jenen Durchmesser verschieden sind. Dieses Verfahren ist sehr zu empfehlen, und man kann dann, wie Bessel thut, den Durchmesser des Feldes durch eine einfache Reihe darstellen. Sobald indessen jener den Declinations-Unterschied der Sterne um mehr als eine Minute übertrifft, gebraucht doch Rec. lieber die Formeln \[
\frac{a'-a}{2d}=\mathrm{tang}\,\varphi,\qquad \frac{a'+a}{2d}=\mathrm{tang}\,\psi,\qquad \frac{d}{\cos\varphi\cos\psi}=D
\] wo $a$, $a'$ die beiden Chorden, $d$ den Declinations-Unterschied der Sterne, $D$ den Durchmesser des Gesichtsfeldes bedeutet. \textsc{Bessel}'s eigene Beobachtungen gehen vom 4. October 1807 bis zum 24. Februar \SourceLeaf{536}1808, und es ist bei ihrer Mittheilung die sehr lobensverthe Manier gebraucht, so viel von ihrem Detail beizufügen, dass man sie künftig bei genauerer Kenntniss der verglichenen und grössten Theils aus der Histoire céleste entlehnten Sterne einer neuen Reduction unterwerfen kann. Hierauf folgen die Beobachtungen des Dr. \textsc{Olbers} in Bremen, die vom 8. October 1807 bis zum 14. Februar 1808 gehen; \textsc{Thulis} Beobachtungen in Marseille vom 22. September bis 2. October; \textsc{Wisniewsky}'s Beobachtungen in Petersburg vom 18. bis 27. März, jene sowohl als diese, vom Verfasser neu reducirt, und endlich, um die im Februar 1808 noch zurückgebliebene Lücke auszufüllen, \textsc{Orianis} Beobachtungen in Mailand vom 13. zum 28. Februar.

Der \emph{zweite Abschnitt} ist der Bestimmung der wahren Bahn des Cometen gewidmet, und zerfällt wieder in drei Abtheilungen. Die \emph{erste Abtheilung} gibt parabolische und rein elliptische Elemente des Cometen. Die ersten parabolischen Elemente berechnete der Verfasser schon am 23. October; die Fortsetzung der Beobachtungen erlaubte am 6. November schon eine Verbesserung derselben, und im März 1808 berechnete er ein drittes System von parabolischen Elementen, die sich an die ersten und letzten ihm damals bekannten Beobachtungen, und an die Rectascensionen der mittleren, möglichst genau anschlossen. Bei den mittleren Declinationen aber blieb noch ein Unterschied von Einer Minute zurück, der sich zwar unter die fünf übrigen Stücke noch sehr würde haben vertheilen lassen, aber doch zu gross geblieben sein würde, um nicht daraus auf das Vorhandensein einer merklichen Abweichung von der Parabel schliessen zu müssen. \textsc{Bessel} liess demnach die parabolische Hypothese fahren und bestimmte, unabhängig von derselben, ein viertes System von Elementen, die eine Ellipse von 1953 Jahren Umlaufszeit ergeben, und deren Vergleichung mit allen Lilienthalschen und Bremischen Beobachtungen nirgends andere Unterschiede bemerken liessen, als so unregelmässig laufende, dass sie nur den Fehlern der einzelnen Beobachtungen zur Last gelegt werden konnten. Bald nachher zeigten doch die ersten Marseiller Beobachtungen, dass jene Elemente die Rectascensionen etwas zu gross gaben, wodurch er zur Berechnung eines fünften Systems von Elementen veranlasst wurde: die Umlaufszeit wurde dadurch auf 1483 Jahre vermindert. Diese Elemente stellten alle Beobachtungen, 91 an der Zahl, so gut dar, dass sich nirgends mehr mit Sicherheit ein entschiedener Unterschied angeben liess, und so war hierdurch die Bahn so genau bestimmt, als es durch rein elliptische Elemente möglich war.

Indessen war es zu erwarten, dass die Berücksichtigung der Störungen, welche der Comet während seiner Sichtbarkeit von den Planeten unsers Sonnensystems erleiden musste, eine um so beträchtlichere Modification dieser Resultate hervorbringen würde, weil die Umlaufszeit so stark von dem kleinen Unterschiede der Excentricität von der Einheit abhängt, und eine auch an sich kleine Aenderung dieses Unterschiedes die Umlaufszeit enorm ändern kann. Es war daher eine interessante Arbeit, diese Störungen gehörig mit in Rechnung zu nehmen, weil sich sonst über die Grenzen, innerhalb welcher man die Umlaufszeit einschliessen dürfe, durchaus gar nicht urtheilen liess. Auf diesen Zweck beziehen sich die beiden folgenden Abtheilungen. In der \emph{zweiten Abtheilung} nemlich entwickelt der Verfasser im Allgemeinen die Methode, die Störungen des Cometen zu berechnen. Bekanntlich ist es hierbei am zweckmässigsten, die Bahn des Cometen als eine veränderliche Ellipse zu behandeln, die täglichen Aenderungen der einzelnen Bestimmungsstücke durch die störenden Kräfte numerisch zu berechnen, und das Collect derselben durch mechanische Quadraturen zu bestimmen. Die Formeln für die augenblickliche Aenderung der Elemente einer elliptischen Bahn durch störende Kräfte sind zwar bereits von mehreren Geometern entwickelt: indess wird man die neue Entwickelung, die \textsc{Bessel} in dieser Abtheilung mit grosser analytischer Eleganz gibt, um so mehr mit Vergnügen lesen, da bei einer Ellipse von so grosser Excentricität zum Theil eine Abänderung nöthig war. Statt der Epoche der mittlern Länge oder Anomalie wurde nemlich die Durchgangszeit durch das Perihelium gewählt, und der sich leicht erge\SourceLeaf{537}bende Ausdruck für die augenblickliche Störung derselben, welcher durch obigen Umstand zur Rechnung nicht gut angewandt werden konnte, bedurfte erst einer besondern Umformung.

In der \emph{dritten Abtheilung} macht endlich der Verfasser die Anwendung dieser Methode auf den Cometen von 1807, indem er zuerst die Summe der störenden Kräfte für sieben gleich weit von einander abliegende Zeitpunkte vom 22.\ September 1807 bis 20.\ März 1808 berechnet. Er sagt nicht ausdrücklich, ob er alle, oder wie viele der Planeten er dabei berücksichtigt habe; Rec.\ hätte gewünscht, den Betrag jedes einzelnen Planeten dabei besonders erwähnt zu finden, da eine Uebersicht dieses Verhältnisses nicht uninteressant gewesen wäre. Hierauf berechnet er für jene sieben Zeitpunkte den Betrag der täglichen Aenderung der einzelnen sechs Elemente. Indem er nun sein zuletzt gefundenes System von Elementen als für den 22.\ September 1807 gültig ansah, berechnete er, nach vorgängiger mechanischer Integration der augenblicklichen Aenderungen derselben, ihre Werthe für sechs verschiedene Zeitmomente vom 28.\ September 1807 bis 23.\ März 1808, für welche er zugleich die geocentrischen Oerter jedesmal aus allen benachbarten beobachteten mit grösster Sorgfalt bestimmt hatte. Die Vergleichung dieser Oerter mit der Rechnung nach jenen veränderlichen Elementen, nebst den numerisch entwickelten Differentialverhältnissen zwischen den einzelnen Elementen und den berechneten geocentrischen Längen und Breiten, gab ihm sodann 12 linearische Gleichungen zwischen den noch zu bestimmenden Correctionen der Elemente, woraus die Werthe dieser Correctionen, mit gebührender Rücksicht auf den comparativen Werth der einzelnen Positionen, nach der Methode der kleinsten Quadrate abgeleitet wurden. So fanden sich endlich die Definitiv-Elemente für den 22.\ September, wovon wir hier nur die halbe grosse Axe $= 143{,}195$, und die Umlaufszeit $= 1713\tfrac{1}{2}$ Jahre, anführen. Diese Umlaufszeit wird indessen durch die fortwirkenden Störungen immerfort abgeändert, so dass sich, wenn auch übrigens dieses Resultat vollkommene Zuverlässigkeit hätte, doch noch ohne eine endlos weitläuftige Rechnung die Zeit der wirklichen Wiederkehr nicht würde angeben lassen. Nach einem Ueberschlage der Störungen durch den Jupiter in den Jahren 1808..1815, wogegen die Störungen durch die andern Planeten unmerklich werden, findet \textsc{Bessel} schon eine Verkürzung der Umlaufszeit von 170 Jahren.

Zum Schluss stellt der Verfasser noch eine Untersuchung über den Grad der Genauigkeit, welchen man diesen Resultaten beilegen darf. Wenn man bei den drei ersten Normalörtern am 28.\ September, 22.\ October und 11.\ November einen Fehler von $5''$ bei der Länge und Breite bei den Oertern vom 8.\ December und 21.\ Februar einen doppelt, und bei dem Orte vom 23.\ März einen vier Mal so grossen voraussetzt, so könnte bei der ungünstigsten Combination derselben die Umlaufszeit auf 2157 Jahre vermehrt, oder auf 1404 Jahre vermindert werden; eine solche Combination ist aber nur möglich, allein höchst unwahrscheinlich. Auch ist, da \textsc{Bessel} so zahlreiche Beobachtungen mit solcher Sorgfalt zur Bestimmung der Normalörter benutzte, kaum zu fürchten, dass noch grössere Fehler bei denselben, als die eben angenommenen, zurückgeblieben sein sollten. Und so ist man wohl befugt, zu hoffen, dass die vorhin bestimmte Umlaufszeit auf hundert Jahre genau sein, und der Comet ungefähr im vierunddreissigsten Jahrhundert unserer Zeitrechnung wiederkehren werde. Dankbar werden sich dann unsere späten Nachkommen an \textsc{Bessel}'s Arbeit erinnern.

Diese letztern Betrachtungen führen nun zugleich auf eine leichte Art zu der vollkommenen Ueberzeugung, dass der Comet \emph{gewiss} eine Ellipse, und keine Parabel oder Hyperbel beschreibt; um eine nicht elliptische Bahn für möglich zu halten, müsste man \emph{wenigstens} siebenfach grössere Fehler bei den Normalörtern zulassen, und dies kann man ohne Bedenken geradezu für unmöglich erklären. So gehört also (bei der jetzigen Vollkommenheit der \emph{practischen Astronomie}) der Comet von 1807 zu den wenigen, von denen wir mit Gewissheit behaupten können, dass, wenn nicht andere Kräfte als die, welche wir kennen, auf sie einwirken, sie dereinst gewiss wiederkehren werden.
